% !TEX root = ../nc-quadric-surface.tex
\subsection{二次被覆の有限判定}
\label{subsec:finite-ambient-test}

\begin{theorem}
\label{thm:finite-ambient-test}
$B=T(V)/(K)$ を左右 noetherian な連結二次 Koszul 環とし、
$H_B(z)=(1+z)/(1-z)^3$ とする。
超平面 $J\subset K$ に対して $A=T(V)/(J)$ と置き、
$K/J$ の非零元を $g$ とする。次は同値である。
\begin{enumerate}
\item $A$ は型 $(1,4,6,4,1)$ の環であり、$g$ は左右正則正規である。
\item 次の有限次元の条件が成り立つ。
\begin{enumerate}
\item $A_3$ で $gV=Vg$ である。
\item $A_i\to A_{i+2}$ の $g$ の左右乗法は $i=1,2$ で単射である。
\item 二次双対 $A^!=T(V^*)/(J^\perp)$ の次数零から五までの次元は
$1,4,6,4,1,0$ である。
\item $A^!_i\otimes A^!_{4-i}\to A^!_4$ の積の組合せが $i=1,2,3$ で非退化である。
\end{enumerate}
\end{enumerate}
従って $B$ の AS Gorenstein 性を先に仮定せず、適合被覆の存在は、
$\operatorname{Gr}(6,K)$ 上で有限個の階数条件を満たす点の有無で判定できる。
\end{theorem}

\begin{proof}
必要性は正規性と Koszul 双対の Hilbert 級数および Frobenius 性から従う。
逆向きに、一次生成性から $gA=Ag$ である。
$\operatorname{Tor}$ の環変換スペクトル系列を用いる。
右 $A$-加群 $B$ に対し $A(-2)\xrightarrow{g\cdot}A\to B\to0$ は極小表示である。
左乗法に核があると仮定し、その最小次数を $d$（$A$ の元の次数）とする。
$g\ne0$ と単射の条件から $d\ge3$ である。
極小分解の第二項から、$\operatorname{Tor}^A_2(B,k)$ の
最小非零内部次数は $d+2$ となる。
また $\operatorname{Tor}^A_1(B,k)\simeq k(-2)$ は
内部次数二だけに集中する左 $B$-加群である。

\[
 E^2_{p,q}=\operatorname{Tor}^B_p(k,\operatorname{Tor}^A_q(B,k))
 \ \Longrightarrow\ \operatorname{Tor}^A_{p+q}(k,k)
\]
において、$E^2_{0,2}$ の内部次数 $d+2$ は非零である。
実際、この最小次数には $B_{\ge1}$ による積が入らない。
この部分に届く可能性のある微分は $(2,1)$ および $(3,0)$ からであるが、
$B$ が Koszul なのでそれぞれ内部次数四および三に集中する。
$(0,2)$ から出る微分はなく、内部次数 $d+2\ge5$ の元は残る。
これは二次環 $A$ の $\operatorname{Tor}^A_2(k,k)$ が
内部次数二に集中することに反する。
反対環にも同じ議論を適用して、$g$ の左右正則性を得る。

従って $B$ の右 $A$-自由分解は長さ一であり、環変換スペクトル系列には
$q=0,1$ の二行のみが残る。$b_p=\dim(B^!)_p$ とすると、
$E^2_{p,0}$ は内部次数 $p$ に $b_p$ 次元、
$E^2_{p,1}$ は内部次数 $p+2$ に $b_p$ 次元ある。
全次数 $n$ の対角内部次数 $n$ に寄与するのは
$E^\infty_{n,0}=\ker(d_2:E^2_{n,0}\to E^2_{n-2,1})$ のみである。
二次環の対角 Tor の次元は二次双対の次元と一致する。
$A^!$ は一次生成なので、次数五の消滅からそれより上も消滅する。
よって仮定により
\[
 \dim E^\infty_{n,0}
 = [z^n](1+z)^4
 = b_n-b_{n-2},
\]
ただし負添字の $b_p$ は零とする。
最後の等号は $H_{B^!}(z)=(1+z)^3/(1-z)$ から従う。
従って全ての $d_2:E^2_{n,0}\to E^2_{n-2,1}$ は全射であり、
非対角の Tor は残らない。これで $A$ は Koszul となる。
同じ対角次元を持つ反対環にも適用できる。

正則正規元で割った環 $B$ が左右 noetherian なので、$A$ も左右 noetherian である。
実際、$(g)$ 進次数環は $B$ 上の skew polynomial 環であり、
斉次イデアルでは固定内部次数でフィルトレーションが有限なので
初期イデアルから有限個の生成元を持ち上げられる。
一般のイデアルには内部次数による最高次成分を用いる。
二次双対は最高次数四で積の組合せが完全なので graded Frobenius である。
Koszul 複体を双対化すると、この組合せによって左右の複体が
内部次数四と複体次数四の移動を除いて互いに同定される。
従って $\Ext_A^j(k,A)$ および反対環の Ext は、$j=4$ の $k(4)$ を除き零となる。
Koszul 性と二次双対の次元により自明加群の極小分解は
長さ四で所要の型となる。任意の有限生成次数付き加群の極小自由分解の長さも、
反対側の $k$ の分解を用いて Tor を計算すれば四以下である。
非次数付き有限生成加群には生成元からの次数フィルトレーションを入れ、
付随次数加群の自由分解を持ち上げると同じ上界を得る。
従って巡回加群に対する射影次元の上界から大域次元も四となる。
正則二次切断から $H_A(z)=(1-z)^{-4}$ が従う。
\end{proof}
